\subsection{Sprint 4}

Na Sprint 4, precisei assumir mais responsabilidades do que eu esperava. Acabei fazendo o deploy do backend para a apresentação, porque os AGES III não respondiam ou alegavam que estavam ocupados. Como a entrega precisava acontecer, não dava para simplesmente esperar indefinidamente. Assumi essa frente para garantir que o projeto tivesse condições de ser apresentado e para evitar que uma pendência crítica comprometesse todo o trabalho acumulado ao longo do semestre.

Além do deploy, desenvolvi de ponta a ponta a autenticação usando um modelo básico com JWT. Também criei o CRUD de arquitetos e finalizei grandes pendências relacionadas à autenticação, incluindo ajustes finais necessários para que o fluxo funcionasse de maneira mais consistente. Essas tarefas exigiram bastante atenção, porque estavam diretamente ligadas à parte administrativa do sistema e ao controle de acesso. Mesmo não sendo ideal concentrar tantas entregas em uma pessoa na última sprint, foi o caminho possível para impedir que o projeto ficasse incompleto.

O ponto mais crítico foi a falta de interesse de alguns colegas, algo que já vinha acontecendo, mas que ficou ainda mais evidente nessa sprint. Quem já demonstrava desinteresse praticamente mostrou zero envolvimento nessa etapa final. O maior impacto veio do sumiço dos AGES III por um bom tempo, o que me deixou várias vezes falando sozinho e sem retorno claro sobre as tarefas. Sem comunicação, acabou sobrando para os AGES IV, e eu terminei fazendo o que estava faltando, que não foi pouca coisa.

Essa situação foi frustrante porque, em uma sprint final, o esperado seria que o time estivesse mais presente e comprometido com a entrega. Em vez disso, precisei lidar com pendências técnicas, ausência de resposta e responsabilidades que deveriam ter sido melhor distribuídas. Ainda assim, preferi resolver o que fosse possível em vez de deixar o projeto chegar incompleto na apresentação.

Em relação às lições aprendidas, sinceramente não sinto que aprendi algo novo nessa sprint. Foi mais uma etapa de executar sob pressão e corrigir problemas que não deveriam ter chegado tão acumulados ao final. Acredito que nosso trabalho como AGES IV poderia ter sido melhor, mas combater a falta de interesse de colegas não é algo simples e muitas vezes não é previsível.
